% Contre-drone
\section{Contre-drone : détecter des drones}
\label{sec:contre-drone}

Sur les 17 articles de contre-drone du corpus, 11 relèvent du sol $\rightarrow$ air ou
ne déclarent pas le capteur. Les deux surveys du domaine
\citep{arxiv-2504.11967-securing-skies-comprehensive-survey-anti,arxiv-2507.10006-vision-anti-unmanned-aerial-technology}\pp{}
ne font pas la distinction : ils mêlent Det-Fly et ARD-MAV, pris depuis un drone, aux jeux
pris du sol, et ne discutent aucune plateforme embarquée chiffrée. Cette section suit donc
le critère du profil de la base : seul l'air $\rightarrow$ air relève de l'IA embarquée.

\subsection{Sol $\rightarrow$ air : des références de contraste}

Les détecteurs sol $\rightarrow$ air du corpus sont tous mesurés sur GPU de bureau ou sur
ordinateur portable :
\begin{itemize}
  \item YOLOv4 élagué \citep{2021-liu-real-time-small-drones-detection} : 90,5 mAP50 et
    69\,FPS sur Tesla P100, sur un jeu propre de deux drones DJI \src{Table~4} ; l'élagage
    fait chuter le rappel sur les petits drones à 30,7\,\%, qu'une augmentation ciblée
    remonte à 72,6\,\% \src{§4.3} ;
  \item F-YOLOv7 \citep{2025-du-f-yolov7-fast-robust-real} : 98,8 d'AP50 de validation sur
    DUT Anti-UAV avec soustraction de fond, sur RTX 4090 \src{Table~2} ;
  \item DotD-YOLOv9-C \citep{2024-gao-dataset-detection-method-unmanned-aerial} : 92,6 de
    mAP@0.5 sur Anti2, sur RTX 4060 portable, sans aucun temps d'inférence \src{Table~5} ;
  \item YOLO-GCOF \citep{2025-yu-yolo-gcof-lightweight-low-altitude} : 91,5 de mAP@50 sur
    DUT Anti-UAV\med{}, annoncé pour Raspberry Pi 4B sans aucun chiffre sur la carte ;
  \item soustraction de fond + MobileNetV2 \citep{2020-seidaliyeva-real-time-accurate-drone-detection} :
    F1 0,742 sur Drone-vs-Bird, 9\,FPS en 1080p sur un matériel non nommé \src{Table~1, §3.4} ;
  \item réseau de caméras stéréo panoramiques
    \citep{2023-zheng-detection-localization-tracking-multiple-mavs} : 9\,FPS avec une RTX
    3070 par groupe de 8 objectifs \src{§VII-B.3}\med ;
  \item fusion IR + visible + audio \citep{arxiv-2007.07396-real-time-drone-detection-tracking}
    sur portable : la fusion à deux capteurs supprime les fausses alarmes relevées en
    10 minutes sans drone \src{§V-D}.
\end{itemize}
Le seul portage FPGA est celui de \citet{2024-nguyen-fpga-soc-implementation-yolov4-flying}
sur ZCU104 (section~\ref{sec:fpga}), sur un jeu d'images d'Internet et de caméras au sol.
Trois idées de ces travaux restent utiles pour l'embarqué : l'augmentation « petits
objets » après élagage, l'appariement DotD sans coût d'inférence, et les jeux
multi-classes (drone, oiseau, avion) pour mesurer les fausses alarmes. En revanche, la
soustraction de fond, centrale dans trois d'entre eux, suppose une caméra fixe : elle ne
fonctionne pas depuis un drone en mouvement, ce que reconnaissent
\citet[§2.1]{2020-seidaliyeva-real-time-accurate-drone-detection}.

\subsection{Air $\rightarrow$ air : le cas embarqué}

\paragraph{Évaluer les détecteurs.} \citet{2021-zheng-air-air-visual-detection-micro}
posent explicitement la distinction : en air $\rightarrow$ air, le fond peut être urbain
ou naturel et la cible faire moins de 10$\times$10 pixels \src{§I}. Sur Det-Fly, les huit
détecteurs évalués vont de 69,5 (RefineDet) à 82,4 d'AP (Grid R-CNN), YOLOv3 étant à
72,3 d'AP pour 32\,ms sur RTX 2080Ti \src{§V-A, §V-B}\med. L'AP moyenne des huit passe de
88,3\,\% sur fond de ciel à 62,0\,\% sur fond urbain \src{Table~III}. Aucun détecteur n'est
porté à bord, alors que le choix de ResNet-50 est motivé par l'embarqué \src{§IV}.

\paragraph{Exploiter le mouvement.} En air $\rightarrow$ air, les méthodes les plus
précises combinent apparence et mouvement :
\begin{itemize}
  \item Dogfight \citep{arxiv-2103.17242-dogfight-detecting-drones-drones-videos}\pp{}
    (CVPR~2021 probable) segmente les drones image par image, puis retrouve des candidats
    par frontières de mouvement et les valide par un I3D sur des cuboïdes de 8 images. Le
    F1 passe de 0,90 (apparence) à 0,92 (apparence et mouvement) sur NPS-Drones
    \src{Table~3}. Aucun temps de calcul n'est rapporté ; la fenêtre temporelle impose au
    moins 4 images de latence \deduit.
  \item GLAD \citep{arxiv-2312.11008-global-local-mav-detection-challenging}\pp{} combine
    un YOLOv5s global, un module de mouvement avec compensation de l'ego-motion par
    homographie, et un YOLOv5s local sur une région de 300$\times$300 suivie par Kalman. Sur
    ARD-MAV, il obtient 0,80 d'AP50 contre 0,56 pour YOLOv5s à 640 et 0,73 pour TPH-YOLOv5l
    à 1536 (12,8\,FPS) \src{Table~IV}. Sur Jetson Xavier NX, en TensorRT, il tourne à
    23,6\,FPS en moyenne, mais à 5,1\,FPS quand le module de mouvement global est appelé
    \src{Table~VIII} ; la précision sur la carte et la puissance ne sont pas mesurées.
  \item GL-YOMO \citep{arxiv-2411.02582-real-time-detection-small-uavs}\pp{} suit la même
    logique global/local avec flot optique et Kalman, à 21,6\,FPS sur Xavier NX
    \src{§IV-E}. Sur Drone-vs-Bird, son AP (0,763) est sous celle de RT-DETR (0,773), à un
    seuil d'IoU de 0,1 \src{Table~II}. Le point de vue du capteur n'est pas déclaré.
\end{itemize}
GLAD est le seul travail air $\rightarrow$ air du corpus avec un chiffre embarqué. Sa
décomposition montre que la partie non neuronale (homographie, Lucas-Kanade, morphologie)
devient le goulot sur Jetson, ce qui plaide pour son accélération matérielle \deduit.
L'idée global $\rightarrow$ local contourne aussi la perte de résolution du
sous-échantillonnage à 640\,px sans tête P2 : un petit YOLO sur une région de
320$\times$320 est un candidat naturel pour un FPGA \deduit.

\subsection{Bilan}

Aucun détecteur air $\rightarrow$ air du corpus n'est mesuré sur SoC FPGA ou MCU, et
aucun n'a de puissance mesurée. Aucun n'est évalué en boucle fermée sur un intercepteur.
Il manque aussi un jeu air $\rightarrow$ air multi-classes (drones, oiseaux, avions) qui
permette de mesurer les fausses alarmes d'un intercepteur ; Anti2
\citep{2024-gao-dataset-detection-method-unmanned-aerial} a les classes mais mêle les
points de vue. C'est le manque le plus net de la base (section~\ref{sec:gaps}).
